% Retained technical source supplement; not input by the focused main.
% Complete readable derivations are Supplementary Notes S4 and S5 in docs.

\subsection{Inherited surface source and the scope of its reduction}
The inherited BHSM source organization assigns the surface and gravitational
contributions to explicit slots:
\begin{align}
 \mathcal L_\psi\psi&=J_{\rm grav}+\mathcal E_\Sigma[Q_\Sigma]+\xi,\nonumber\\
 Q_\Sigma&=\frac{c^2}{8\pi}\int_{\Sigma_\star}
 H_\Sigma\Delta K_{uu}\,\dd A,\nonumber\\
 h_{\rm topo}&=-\mathcal A_g^{-1}\mathcal B_{g\psi}\psi .
 \label{eq:inherited-bhsm-source}
\end{align}
These are the retained source and constrained metric-response definitions,
not a newly closed constitutive law. The surface functional uses the seam
curvature data; its export operator, orientation and full action
normalization remain to be matched to the local field variables. Here
$\psi$ names that inherited topographic variable, not the independent
rolling scalar $T$, the target field $\eta$, or the metric coordinate
$\psi_{\rm geo}$ introduced below. The local metric Hessian in this source
organization is also distinct from the cosmological ADM constraint block.
No stochastic forcing or population covariance is required for the
specified deterministic encounter. The term $\xi$ remains a source slot;
it is not supplied with an invented realization or state law.

\begin{table}[htbp]
 \centering
 \caption{Status of the source-to-observable construction.}
 \begin{tabular}{>{\raggedright\arraybackslash}p{0.22\linewidth}>{\raggedright\arraybackslash}p{0.69\linewidth}}
 \toprule
 Status & Statement \\
 \midrule
 Adopted BHSM premise & Environmental light traverses an envelope event;
 conventional and scalar descriptions represent its surface contributions.\\
 Derived effective result & The declared R1 action fixes the local matter source,
 regular scalar branch and metric response, hence the optical transfer.\\
 Derived equivalence & A positive effective mass reproduces that retained
 metric and optical field.\\
 Derived source condition & Cap momentum and Schur forcing specify the
 required pullback; positive same-domain lifts prove its microscopic
 completion is not unique from the leading coefficients.\\
 Open common-source map & Reduction of the BHSM surface and boundary data to
 the normalized R1 source, with the same independent variables and validity range.\\
 \bottomrule
 \end{tabular}
\end{table}

The calculation below executes the effective problem and its optical
application. The derived common-source conditions specify the physical
selection data consumed by a further application within the BHSM program.


\section{Exact Bulk Continuation, Environment and Matching}
\label{app:continuation}
The retained action also specifies the next field problem. The field
equations below use $c=\hbar=1$, restoring $c$ in scale comparisons. Its
exact bulk equations can be derived without supplying the unfinished microscopic
surface-to-R1 map. They are distinct from an evolved physical solution
with a specified material law and global state. The full spherical
component derivation is supplied in
\path{docs/strong_lensing_r1_screened_20261020/R1_REMAINING_FIELD_DERIVATION_20261020.md}.

\subsection{The current balance and the rolling reference}
Put $U_\mu=\nabla_\mu T$, $Y=U_\mu U^\mu$,
$\Pi_{\mu\nu}=\nabla_\mu\nabla_\nu T$, $S=\Box T$,
$\mathcal K_\Pi=\Pi_{\mu\nu}U^\mu U^\nu$ and $g_3=2/\lambda$.
The exact metric equation and cubic stress are
\begin{align}
 \Mpl^2G_{\mu\nu}&=T^m_{\mu\nu}+U_\mu U_\nu
             -g_{\mu\nu}(Y/2+V)+T^{(3)}_{\mu\nu},\nonumber\\
 T^{(3)}_{\mu\nu}&=g_3[-S U_\mu U_\nu-g_{\mu\nu}\mathcal K_\Pi
       +2U_{(\mu}\Pi_{\nu)\alpha}U^\alpha].
 \label{eq:continuation-einstein}
\end{align}
Equation~\eqref{eq:scalar-covariant} has the exact current form
\begin{equation}
 J^\mu=U^\mu+g_3(\Pi^{\mu\nu}U_\nu-SU^\mu),\qquad
 \nabla_\mu J^\mu=-\mu^3.
 \label{eq:continuation-current}
\end{equation}
On the homogeneous reference $J^0=-J$ for the positive current in
Equation~\eqref{eq:current}, reproducing its source sign.
For a spherical patch
$ds^2=-N^2dt^2+A^2dr^2+r^2d\Omega^2$, the actual balance is
\begin{equation}
 \partial_t(NAr^2J^t)+\partial_r(NAr^2J^r)=-\mu^3NAr^2.
 \label{eq:continuation-balance}
\end{equation}
An exactly static metric with $T=vt+\pi(r)$ and constant $v$ would remove
the temporal term. A regular center would then require
\begin{equation}
 J^r(r)=-\frac{\mu^3}{NAr^2}\int_0^rN(s)A(s)s^2\,\dd s,
 \label{eq:continuation-static}
\end{equation}
which grows as $-\mu^3r/3$ in a static asymptotically flat region, rather
than giving the decaying isolated excess. Moreover $V_0-\mu^3vt-\mu^3\pi$
makes the stress time dependent for nonzero $\mu^3v$. The rolling
background time-current balance must therefore be retained before its
homogeneous contribution is subtracted. The executed quasi-static
solution does this at its retained order; appending a static GR2 potential
does not solve this full continuation. Small $H_lr_V/c$ supports the
local freezing approximation but is not a physical delay-error bound.

\subsection{The environmental contraction}
At the retained quasi-static order let $\pi_T=\pi_L+\pi_E$ and
$\delta\rho=\rho_L+\rho_E$. Subtracting the environmental equation from
the total scalar equation gives the actual local application
\begin{align}
 C\Delta\pi_L
 &-\frac2\lambda[(\Delta\pi_L)^2-\Pi_L:\Pi_L]\nonumber\\
 &-\frac4\lambda[(\Delta\pi_E)(\Delta\pi_L)-\Pi_E:\Pi_L]
   =-\frac{b_T\rho_L}{2\Mpl^2},
 \label{eq:continuation-environment}
\end{align}
where $\Pi_{L,ij}=\partial_i\partial_j\pi_L$ and similarly for $E$.
The environmental input is the full spatial Hessian in
$[(\operatorname{tr}\Pi_E)\delta^{ij}-\Pi_E^{ij}]
\partial_i\partial_j\pi_L$, not its trace alone. Constant or
uniform-gradient scalar data have no such cross term, whereas a tidal
Hessian generally does. A projected optical tide does not determine this
three-dimensional contraction. Physical environmental density and
boundary/Cauchy data must be solved with the same scalar and metric
constraints; their total principal operator must also be checked.
The optical nuisance tides in the main text do not perform that solve.

\subsection{Material and outer-domain defining data}
An evolved perfect fluid consumes a material relation, for example
$\rho_m=\rho_m(n)$ and $p_m=n\,\dd\rho_m/\dd n-\rho_m$, together with
$\nabla_\mu T_m^{\mu\nu}=0$. At static equilibrium the radial application
is $p_m'=-(\rho_m+p_m)\nu'$ for $N=e^\nu$.
A common fixed barotropic law with $\delta\rho_m=0$ would give
$\delta p_m=0$, whereas Equation~\eqref{eq:pressure} gives a positive
pressure change. The retained comparison therefore prescribes common
density, not a common fixed material equation of state. A physical
material continuation must allow density readjustment, specify additional
support, or explicitly retain the prescribed-stress comparison.

In an overlap worldtube $r_V\ll r_m\ll c/H_l$, the regular branch fixes
its far-zone derivative data without a new local charge:
\begin{align}
 \pi_T'(r_m)&=-\frac{b_TGM}{Cr_m^2}
       [1+O((r_V/r_m)^3)],\nonumber\\
 \Phi'(r_m)=\Psi'(r_m)&=\frac{GM}{r_m^2}
       [1+\epsilon_\infty+O((r_V/r_m)^3)].
 \label{eq:continuation-overlap}
\end{align}
Smooth matching also requires the induced metric and extrinsic curvature,
the scalar value and normal derivative, and the current balance in
Equation~\eqref{eq:continuation-balance}. Scalar additive data and clock
normalization belong to the outer state. A physical surface source would
instead require its variational junction conditions; no moving-seam law
is assumed here.

The flat local constraint cannot be extended unchanged to the entire
compact slice. Its formal compact integral would require
$\int\delta\rho\,\dd V=0$, which the isolated positive excess does not
satisfy. The actual outer problem must retain curvature, time dependence
and the homogeneous response, or a declared compensating excess. This
is a compatibility check of the naive extension, not a prohibition of
positive local mass in a closed cosmology. The outer solution must match
both the local field and the source/observer clock protocol. Thus the
exact bulk and environmental equations are available, while their
material and global state data must be specified before a further
physical solution is claimed.

\section{Common-Source Pullback and a Nonuniqueness Construction}
\label{app:source-reduction}
The original cap reduction is a constrained critical value with retained
boundary traces and kernel coordinates. Its actual matcher variation
defines the oriented canonical momenta by
$\delta S_{5,\epsilon}^{\rm on\mbox{-}shell}
=\int\Pi_\epsilon^{ab}\delta\gamma_{\epsilon ab}$ and
$S_{\rm match}=\sum_\epsilon\int\Lambda_\epsilon^{ab}
(h_{ab}-\gamma_{\epsilon ab})$. It gives
\begin{equation}
 \Pi_\epsilon^{ab}=\Lambda_\epsilon^{ab},\qquad
 2C_\partial G^{ab}[h]=T_m^{ab}+\tau_{\rm cap}^{ab},\qquad
 \tau_{\rm cap}^{ab}=\frac2{\sqrt{-h}}\sum_\epsilon\Pi_\epsilon^{ab}.
 \label{eq:cap-source-pullback}
\end{equation}
Thus $u_a u_b\delta\tau_{\rm cap}^{ab}$ supplies the metric energy
source once a cap solution and observer are specified. The cap scalar
normal momentum and the interface Euler variation supply different
source slots; their pullback into canonical $T$ must be established.
The cap sum cannot be replaced with the reflection-even mean without
its factor of two. A cap response absorbed into the local action is
accounted for once.

For original coupled rows $A_P\psi+Bq=J_P$,
$B^\dagger\psi+D_Qq=J_Q$, elimination on an invertible owned complement
gives both the operator and its forcing:
\begin{equation}
 (A_P-BD_{Q,\perp}^{-1}B^\dagger)\psi
 =J_P-BD_{Q,\perp}^{-1}J_Q-Bq_0,
 \qquad q_0\in\ker D_Q.
 \label{eq:source-schur}
\end{equation}
Kernel compatibility equations remain explicit. A reduced Hessian
alone does not determine this source. The physical embedding of $T$,
its kinetic normalization, cubic vertex and source pairing must agree
under the same reduction.

The leading R1 coefficients do not select a unique complement. To prove
this, take two mathematical spectral lifts on the same ball
$L=10\pi\ell_0$, both real fields regular at zero and Dirichlet at $L$.
The normalized radial modes are
$f_n=\sin(n\pi r/L)/(r\sqrt{2\pi L})$, $n\ge1$,
with $q=p\ell_0^2=(n/10)^2$. Define
\begin{equation}
 H(q)=\begin{pmatrix}(C+1)q&g\sqrt q\\g\sqrt q&M_c^2+q\end{pmatrix},
 \quad K_{\rm time}=\operatorname{diag}(D,1),
 \quad (M_c^2,g)=(1,1)\ \hbox{or}\ (4,2).
 \label{eq:complement-witness}
\end{equation}
Here $M_c,g,q$ are dimensionless; $D$ is the temporal coefficient in
Appendix~\ref{app:principal}. The same retained source is used and the
complement is explicitly unforced for this counterexample. Neither lift
is selected as a microscopic galaxy state. Each complement is positive,
and its reduced kernel is
$Cq+g^2q^2/[M_c^2(M_c^2+q)]>0$ for every owned $q>0$.
Both leading coefficients are exactly $C,D$, yet after normalizing by
$C$ their fourth-order coefficients are
$B/\ell_0^2=1/C=0.2359125661$ and
$1/(4C)=0.05897814152$, with expansion scales
$p\ell_0^2\ll1$ and $p\ell_0^2\ll4$, respectively.
The common unit $\ell_0$ has not been assigned a physical galaxy scale.
This constructive nonuniqueness result identifies missing physical
selection data rather than inferring them from an unavailable array.

A further direct identification can be excluded. The inherited
distributed seam law
$\nabla\cdot\mathbf F_{\rm top}=a\Theta(r-r_\star)/r^2$ with
$F^r_{\rm top}(r_\star)=0$ integrates to
$F^r_{\rm top}=a(1/r-r_\star/r^2)$. Its growing integrated charge
cannot be identified by a constant normalization with the isolated R1
flux, which decays as $r^{-2}$. A different finite-domain export or
nonlocal source map remains possible. The full action/source trail,
including the interface Euler relation, is supplied in
\path{docs/strong_lensing_r1_screened_20261020/BHSM_COMMON_SOURCE_DERIVATION_20261020.md}.
It specifies the physical embedding, constitutive interface density and
boundary prescription needed for a unique further application.
